% !TeX root = ../main.tex
% Appendix A3: Cumulant Expansion of the Implied Volatility Surface
% Batch #4.9a — Notation hygiene

\section{Cumulant Expansion of the Implied Volatility Surface}\label{app:cumulants}

Equation \eqref{M-eq:effective-vol} follows from a cumulant expansion of the
log-return characteristic function
\citep{jarrow1982approximate,corrado1996skewness,backus2004accounting}. Let
$\varphi(u) = \mathbb{E}[e^{iu X_T}]$ for $X_T = \ln(S_T/S_0)$, with cumulant
generating function $\psi(u) = \ln \varphi(u)$.

In the quantum framework, the characteristic function is computed from the Wigner
function:
\begin{equation*}
    \varphi(u) = \int e^{iux} p(x, T) \, dx
    = \int e^{iux} W(x, p, T) \, dx \, dp.
\end{equation*}

The cumulant expansion of the implied volatility is obtained by matching the
characteristic function of the quantum model to the BSM characteristic function
and solving for the volatility that equates the option prices. The first few
cumulants of the quantum model are:
\begin{align}
    \kappa_1 &= \mu T, \label{eq:cumulant-1} \\
    \kappa_2 &= \sigma_0^2 T + \hbar_{\mathrm{econ}}^2 \cdot \beta
        + \mathcal{O}(\hbar_{\mathrm{econ}}^4), \label{eq:cumulant-2} \\
    \kappa_3 &= \hbar_{\mathrm{econ}} \cdot \gamma T
        + \mathcal{O}(\hbar_{\mathrm{econ}}^3), \label{eq:cumulant-3} \\
    \kappa_4 &= \mathcal{O}(\hbar_{\mathrm{econ}}^2), \label{eq:cumulant-4}
\end{align}
where $\mu$ is the drift, $\sigma_0^2$ is the classical variance,
$\gamma = \omega([\hat{X},\hat{P}]_+)$ and
$\beta = \omega(\hat{X}^2)$ are the quantum correction coefficients
from \eqref{M-eq:effective-vol}. Note that $\kappa_2^{(\hbar^2)} =
\hbar_{\mathrm{econ}}^2\beta$ is time-independent at this order.
To reconcile with Appendix~\ref{app:moyal} Step~2, where the variance
correction is $2\hbar^2\beta_1 = \hbar^2\beta/\tau$ (via $\beta_1=\beta/(2\tau)$):
since $\kappa_2 = \sigma^2_{\mathrm{imp}}\cdot T$ (variance $\times$ horizon $T=\tau$),
a variance correction $\delta\sigma^2 = \hbar^2\beta/\tau$ gives
$\delta\kappa_2 = \delta\sigma^2\cdot T = \hbar^2\beta$, consistent with
$\kappa_2^{(\hbar^2)} = \hbar^2\beta$ here. The two representations are
therefore equivalent; the $\tau^{-2}$ factor in the smile formula
\eqref{eq:full-smile-correction} arises from the BSM inversion step in
Appendix~\ref{app:moyal} Step~3.

The skewness $\kappa_3 / \kappa_2^{3/2}$ diverges as $\mathcal{O}(T^{-1/2})$
when $T \to 0$. This is a \emph{different observable} from the skew slope
$\partial \sigma_{\mathrm{imp}}/\partial x$ discussed in the main text. The
two quantities are related but distinct:

\begin{itemize}
\item The \textbf{standardized skewness} $\kappa_3/\kappa_2^{3/2}$
  diverges as $\mathcal{O}(T^{-1/2})$ because $\kappa_3 \propto T$ while
  $\kappa_2^{3/2} \propto T^{3/2}$, giving $T/T^{3/2} = T^{-1/2}$.

\item The \textbf{skew slope} $\partial\sigma_{\mathrm{imp}}^2/\partial x$,
  defined as the moneyness derivative of implied variance, is given by
  $\hbar_{\mathrm{econ}}\gamma/\tau$ from~\eqref{eq:full-smile-correction},
  which diverges as $\mathcal{O}(\tau^{-1})$.
\end{itemize}

The $\mathcal{O}(T^{-1})$ divergence of the skew slope---claimed in
Corollary~\ref{M-cor:skew-explosion} and central to the empirical predictions
of the paper---arises directly from differentiating the smile formula with
respect to log-moneyness $x$, not from the cumulant ratio. The cumulant
$\kappa_3$ characterises the overall distributional skewness; the smile
slope characterises the local moneyness sensitivity. Both diverge as
$T \to 0$ but at different rates, which is consistent: the $T^{-1}$ slope
is steeper than the $T^{-1/2}$ skewness, reflecting the non-linear BSM
inversion from log-return distribution to implied volatility surface.
The kurtosis $\kappa_4 / \kappa_2^2$ remains $\mathcal{O}(1)$, distinguishing the quantum model from
jump-diffusion models where the kurtosis diverges as $\mathcal{O}(T^{-1})$.
